\section{Related Work}

\subsection{Cache-conscious data}

%\paragraph{Using generational garbage collection to implement cache-conscious data placement in object-oriented languages~\cite{ccgarbage}}
%
Chilimbi and Larus~\cite{ccgarbage} base on an object-oriented language with a
generational garbage collector, which they extend with a heuristic for
copying objects to the TO space.
%
Their heuristic uses a special-purpose graph data structure, the \emph{object affinity
graph}, to identify when groups of objects are accessed by the program
close together in time.
%
When a given group of objects have high affinity in the object
affinity graph, the collector is more likely to place them close
together in the TO space.
%
As such, a goal of their work and ours is to achieve higher data-access locality
by carefully grouping together objects in the heap.
%
However, a key difference from our work is that their approach bases its placement decisions
on an object-affinity graph that is generated from profiling data, which is
typically collected online by some compiler-inserted instrumentation.
%
The placement decisions made by our approach are based on data collected by static analysis of the program.
%
Such an approach has the advantage of not depending on the output of dynamic profiling, and therefore
avoids the implementation challenges of dynamic profiling.
%
A disadvantage of not using dynamic profiling is that the approach
cannot adapt to changing access patterns that are highly input
specific.
%
We leave open for future work the possibility of getting the best of both approaches.

%% There are two properties of object oriented programming that are utilized. Objects are less than 32 KB and that they are not lightweight use objects. Since the cache block sizes these days are 64 KB this means that the fields of the objects need not be optimized for and the only layout optimization should be object level optimization.
%% The paper introduces a concepts of a object affinity graph which is a graph consisting of vertices as live objects
%% % in the garbage collector
%% and edges encode the temporal affinity between the objects. This graph is generated by hooking the application with a real time profiler that tracks these temporal affinities. During a scavenge phase the copying garbage collector then copies these objects over to the ``TO'' space in an order that increases the cache performance. They use a greedy algorithm that traverses the max weighted edge and copyies the objects over in this greedy fashion. Miss rates are improved by 20\% to 41\% and performace by 18\% to 31\% compared to prior work.


%\paragraph{Cache-Conscious Structure Definition~\cite{CSD}}
%
%The work extends the work Cache-Conscious structure layout by focusing on data structures that are larger than the size of a cache block. Previous work worked best for data structures that are around half the size of a cache block. This work builds upon the prior work by extending it with two more techniques. These are \textbf{structure splitting} and \textbf{field reordering}.
%
%Mapping concurrently accessed elements of a struct in the same cache block provides an implicit prefetch. If these elements do not fit in the same cache block then placing them in non-conflicting cache blocks will reduce cache misses.
%
Chilimbi et al.~\cite{CSD} introduce the idea of hot/cold splitting of a data structure, where elements are categorized as being ``hot'' if accessed frequently and ``cold'' if accessed inferquently.
This information is obtained by profiling the program. Cold fields are placed into a new object via an indirection and hot fields remain unchanged. In their approach, at runtime, there is a cache-concious garbage collector~\cite{ccgarbage} that co-locates the modified object instances.
This paper also suggests placing fields with high temporal affinity within the same cache block. For this they recommend  \texttt{bbcache}, a field recommender for a data structure. \texttt{bbcache} forms a field affinity graph which combines static information about the source location of structure field accesses with dynamic information about the temporal ordering of accesses and their access frequency.
% \textbf{Shortcomings: } Applicable to only the pointer world, does not apply to the packed world like in gibbon where the trees are serialized in memory. Not applicable to functional progams. Does not apply to recursive data structures.

% \paragraph{Cache-Conscious Structure Layout~\cite{CSL}}

%Poor reference locality is one major bottleneck in today's programs. Since programmers often do not program keeping in mind the hierarchical  structure of caches in today's architectures, the memory performance of these programs is sometimes far from ideal.
% This paper targets the problem at the source, i.e, solve the problem of poor reference locality within these programs (pointer based programs in C).
%Their proposed technique is improves on prior work in that it requires less programmer effort, less architectural dependence and less source code modifications.
Chilimbi et al.~\cite{CSL} propose two techniques to solve the problem of poor reference locality.
\begin{description}
\item[ccmorph] This works on tree-like data structures, and it relies on the programmer making a calculated guess about the safety of the operation on the tree-like data structure. It performs two major optimizations: clustering and coloring. Clustering take a the tree like data structure and attempts to pack likely to be accessed elements in the structure within the same cache block. There are various ways to pack a subtree, including clustering $k$ nodes in a subtree together, depth first clustering, etc.
% This depends on the program. This technique does not attempt to suggest the optimal clustering based on the program.
Coloring attempts to map simultaneously accessed data elements to non-conflicting regions of the cache.
\item[ccmalloc] This is a memory allocator similar to \texttt{malloc} which takes an additional parameter that points to an existing data structure element which is likely accessed simultaneously. This requires programmer knowledge and effort in recognizing and then modifying the code with such a data element. \texttt{ccmalloc} tries to allocate the new data element as close to the existing data as possible, with the initial attempt being to allocate in the same cache block. It tries to put likely accessed elements on the same page in an attempt to improve TLB performance.
\end{description}
%Using these techniques they are able to get significant improvements in the run-time of the program. Future work is to get rid of programmer effort and be able to analyze the program statically or profiling techniques.

%\paragraph{SHAPES~\cite{Shapes, Shapes2}}
Franco et al.~\cite{Shapes, Shapes2} suggest
that the layout of a data structure should be defined once at the point of initialization, and all further code that interacts with the structure should be ``layout agnostic''. Ideally, this means that performance improvements involving layout changes can be made without requiring changes to program logic. To achieve this, classes are extended to support different layouts, and types carry layout information---code that operates on objects may be polymorphic over the layout details.

\subsection{Data layout description and binary formats}

Chen et al.~\cite{Dargent} propose a data layout description framework \emph{Dargent}, which allows programmers to specify the data layout of an ADT. It is built on top of the Cogent language~\cite{o2016cogent}, which is a first order polymorphic functional programming language. Dargent targets C code and provides proofs of formal correctness of the compiled C code with respect to the layout descriptions. Rather than having a compiler attempting to determine an efficient layout, their focus is on allowing the programmer to specify a particular layout they want and have confidence in the resulting C code.
% Their work however, is set in a different environment and does not evaluate the data layout transformations in view of performance. In addition, they do not automatically optimize for the best layout when it comes to performance. Their work does not support the ordering of recursive fields in the algebraic data type either.

% \paragraph{PADS ~\cite{PADS, Pads2, Pads3}} Something about PADS/ML here.

% \subsection{Protocols, parsing, and binary formats}

Significant prior work went into generation of verified efficient code for interacting with binary data formats (parsing and validating). For example, EverParse~\cite{EverParse} is a framework for generating verified parsers and formatters for binary data formats, and it has been used to formally verify zero-copy parsers for authenticated message formats. With Narcissus~\cite{Narcissus}, encoders and decoders for binary formats could be verified and extracted, allowing researchers to certify the packet processing for a full internet protocol stack. Other work~\cite{GenericPacketDescriptions} has also explored the automatic generation of verified parsers and pretty printers given a specification of a binary data format, as well as the formal verification of a compiler for a subset of the Procotol Buffer serialization format~\cite{VPBdeleware}.

Back~\cite{DataScript} demonstrates how a domain-specific language for describing binary data formats could be useful for generating validators and for easier scripting and manipulation of the data from a high-level language like Java. \cite{PacketTypes} introduce \emph{Packet Types} for programming with network protocol messages and provide language-level support for features commonly found in protocol formats like variable-sized and optional fields.

%Data Representation Synthesis~
Hawkins et al.~\cite{Hawkins} introduce \relc, a framework for synthesizing low-level C\texttt{++} code from a high-level
relational representation of the code. The user describes and writes code that represents data at a high level as relations. Using a decomposition
of the data that outlines memory representation, \relc synthesizes correct and efficient low-level code.

Baudon et al.~\cite{BitStealingMadeLegal} introduce the \ribbit DSL, which allows programmers to describe the layout of ADTs that are monomorphic and immutable.
\ribbit provides a dual view on ADTs that allows both a high-level description of the ADT that the client code follows and a user-defined memory representation of the ADT
for a fine-grained encoding of the layout. Precise control over memory layout allows \ribbit authors to develop optimization algorithms over the ADTs, such as struct packing, bit stealing, pointer tagging, unboxing, etc.
Although this approach enables improvements to layout of ADTs, it is different from \system's: \ribbit focuses on \emph{manually} defining low-level memory representation of the ADT whereas \system \emph{automatically} optimizes the high-level layout (ordering of fields in the definition ADT) relying on \gibbon for efficient packing of the fields.
While \ribbit invites the programmer to encode their best guess about the optimal layout, \system comes up with such layout by analysing access patterns in the source code.
% This has performance enhancing effects due to locality. % Artem: looks disconnected from the previous sentence

\subsection{Memory layouts}


Early work on specifications of memory layouts was explored in various studies of PADS, a language for describing ad hoc data-file formats~\cite{PADS, Pads2, Pads3}.

Lattner and Adve~\cite{poolallocation} introduce a technique for improving the memory layout of the heap of a given C program.
%
Their approach is to use the results of a custom static analysis to enable \emph{pool allocation} of heap objects.
%
Such automatic pool allocation bears some resemblance to our approach, where we use region-based allocation in tandem with region inference, and thanks to static analysis can group fields of a given struct into the same pool, thereby
improving locality in certain circumstances.
%

Floorplan~\cite{Floorplan} is a declarative language for specifying high-level memory layouts, implemented as a compiler which generates Rust code. The language has forms for specifying sizes, alignments, and other features of chunks of memory in the heap, with the idea that any correct state of the heap can be derived from the Floorplan specification. It was successfully used to eliminate 55 out of 63 unsafe lines of code (all of the unsafe code relating to memory safety) of the immix garbage collector.


% While this paper has been concerned with the format and layout of serialized objects and fields in memory, there has been significant research on the serialized layout of data persisted to disk or sent over the network: on the specification of network message protocols~\cite{PacketTypes}, on the efficient parsing of authenticated message formats and binary data~\cite{EverParse}, and so on. Similarly, prior work has explored the verified  with binary data~\cite{DataScript, Narcissus}.

% \subsection{Uncategorized}

% \paragraph{Shapes~\cite{Shapes, Shapes2}}

% \paragraph{hobbit~\cite{Hobbit, Hobbit2}}

% \paragraph{Contiguity types ~\cite{contiguityTypes}}

% \paragraph{The hardness of cache conscious data placement ~\cite{hardnessLayout}}

% \paragraph{Data layout optimization based on the spatio-temporal model of field access ~\cite{9948361}}

% \paragraph{Recursion Twisting ~\cite{recursionTwisting}}

%\paragraph{Automatic Pool Allocation for Disjoint Data Structures ~\cite{poolallocation}}


%The idea is that it is easier to perform intraprocedural optimization at link time when multiple files are being compiled together. The compiler generates data structure graphs of disjoint data structures for each functions, for data structures that are accessed in a type safe manner, it automatically allocates pools for data structures.

% \paragraph{Detecting Conflicts Between Structure Accesses ~\cite{larusAndHilfinger}}

% \paragraph{An analytical cache model ~\cite{analyticalCacheModel}}

% \paragraph{Principles of optimal page replacement ~\cite{popr}}

% \paragraph{Organizing and maintenance of large ordered indexes ~\cite{Bayer1972OrganizationAM}}

% \paragraph{Enhancing performance in a persistent object store ~\cite{eronique1995enhancing}}
